% Section 7.3, from lectures/L07/appendix/b_exercises.md.

\section{Exercises}
\label{c7:sec:exercises}\label{c7:app:b}

\begin{aside}
\textbf{How to work these.} The Reflection exercises are answered on paper, from the chapter. The
Code exercise is a program for the board: build and flash it the way \secref{c2:app:c} (Microchip
Studio) or \secref{c2:app:d} (Linux) describes, or run it in the simulator of
Appendix~\ref{sim:ch} if you have no board yet: its simulation harness (\secref{sim:sec:harness})
saves the EEPROM to a file, so what the program stores survives a restart there too.

\textbf{How to check your work.} The Reflection exercises, and the question the Code exercise asks
about its design, are answered in Appendix~\ref{app:answers}. The Code exercise has no printed
solution: the board tells you whether the program does what the exercise describes. Try each
exercise before you read its answer.
\end{aside}

% ------------------------------------------------------------------------------------------
\exerciseset{Concepts}
\label{c7:set:1}

\reflexercise{Waiting before, not after}
\label{c7:ex:1}
{\raggedright
Why does \code{eeprom_write_byte()} wait on the \code{EEPE} flag before it starts, but
\code{eeprom_read_byte()} doesn't need to wait for anything after its read completes?\par}

\reflexercise{The four-cycle window}
\label{c7:ex:2}
\code{EEMPE} must be followed by \code{EEPE} within four CPU cycles or the write silently fails.
Why does that specifically mean interrupts need to be disabled for that window, and what real
failure would you expect if an ISR fired in between and took longer than four cycles to return?

% ------------------------------------------------------------------------------------------
\exerciseset{Program}
\label{c7:set:2}

\codeexercise{Fewer EEPROM writes}
\label{c7:ex:3}
\secref{c7:a:4}'s \file{eeprom_driver_demo} writes to EEPROM on every single button press. Given
\secref{c7:a:1}'s write-cycle warning, is that actually a good idea for a device expected to run for
years? Extend the demo to reduce wear: keep a small, non-persisted counter of presses since the
last EEPROM write, and only call \code{eeprom_write()} once it reaches \code{5}, instead of on every
press. \code{main()}'s event handling already separates ``an event happened'' from ``write it
out'', build on that rather than restructuring it.
